\documentclass[11pt]{article}

\usepackage[margin=1in]{geometry}
\usepackage{booktabs}
\usepackage{graphicx}
\usepackage{hyperref}
\usepackage{amsmath}
\usepackage{array}
\usepackage{xcolor}

\hypersetup{
    colorlinks=true,
    linkcolor=blue,
    urlcolor=blue,
    citecolor=blue
}

\newcommand{\placeholder}[1]{\textcolor{gray}{\textit{#1}}}

\title{Video Action Recognition with UCF50 using LRCN}
\author{Kevin Wagner}
\date{\today}

\begin{document}

\maketitle

\begin{abstract}
This report describes an LRCN-based approach for video action recognition on the UCF50 dataset.
The model uses a pretrained ResNet backbone to extract spatial frame features and an LSTM to model temporal information across uniformly sampled video frames.
This draft leaves placeholders for final training logs, classification metrics, confusion matrix analysis, and test-set accuracy after the full experiment is run.
\end{abstract}

\section{Introduction}

Video action recognition requires modeling both visual appearance and motion over time.
In this assignment, the goal is to classify videos from the UCF50 dataset into one of 50 action categories.
The implemented approach follows a Long-term Recurrent Convolutional Network (LRCN) design: a convolutional neural network extracts frame-level representations, and a recurrent neural network aggregates those representations over a short temporal sequence.

\section{Dataset}

The UCF50 dataset contains realistic YouTube videos grouped into 50 human action categories.
The dataset was organized with one directory per class:

\begin{verbatim}
UCF50/
    BaseballPitch/
    Basketball/
    ...
\end{verbatim}

\noindent
\textbf{Dataset split.}
The preprocessed frame dataset was split into training, validation, and test subsets using stratified sampling.
The planned split is shown in Table~\ref{tab:data-split}.

\begin{table}[h]
    \centering
    \caption{Dataset split configuration. Replace counts after preprocessing UCF50.}
    \label{tab:data-split}
    \begin{tabular}{lcc}
        \toprule
        Split & Proportion & Number of Videos \\
        \midrule
        Training & 75\% & \placeholder{insert count} \\
        Validation & 10\% & \placeholder{insert count} \\
        Test & 15\% & \placeholder{insert count} \\
        \bottomrule
    \end{tabular}
\end{table}

\section{Preprocessing}

Each video was converted into a fixed-length sequence of image frames before training.
The preprocessing script uniformly samples frames across the full duration of each video and stores them as JPEG images.
For the main experiment, 16 frames per video are used.

\begin{verbatim}
python preprocess.py --input_dir UCF50 --output_dir UCF50_frames --frames_per_video 16
\end{verbatim}

\noindent
The resulting directory structure is:

\begin{verbatim}
UCF50_frames/
    BaseballPitch/
        video_name/
            frame_0000.jpg
            frame_0001.jpg
            ...
\end{verbatim}

During training, frames are resized to $224 \times 224$, randomly flipped, randomly translated, converted to tensors, and normalized using ImageNet statistics.
Validation and test transforms use resizing, tensor conversion, and normalization without random augmentation.

\section{Model Architecture}

The model is an LRCN classifier with two main components:

\begin{itemize}
    \item \textbf{CNN backbone:} A ResNet model extracts one feature vector from each sampled frame.
    \item \textbf{Temporal model:} An LSTM processes the sequence of frame features.
    \item \textbf{Classifier:} A dropout layer and fully connected layer map the final LSTM output to 50 class logits.
\end{itemize}

\noindent
The planned model configuration is shown in Table~\ref{tab:model-config}.

\begin{table}[h]
    \centering
    \caption{Model configuration for the main experiment.}
    \label{tab:model-config}
    \begin{tabular}{ll}
        \toprule
        Component & Configuration \\
        \midrule
        Model type & LRCN \\
        CNN backbone & ResNet-34 \\
        CNN initialization & ImageNet pretrained weights \\
        Recurrent layer & LSTM \\
        LSTM hidden size & 100 \\
        Number of LSTM layers & 1 \\
        Dropout & 0.1 \\
        Number of classes & 50 \\
        Frames per video & 16 \\
        \bottomrule
    \end{tabular}
\end{table}

\section{Training Procedure}

The model is trained using cross-entropy loss and the Adam optimizer.
A ReduceLROnPlateau scheduler lowers the learning rate when validation loss stops improving.
The best model weights are selected using validation accuracy and saved for final test evaluation.

\begin{verbatim}
python run.py --frame_dir UCF50_frames --train_size 0.75 --test_size 0.15 \
  --model_type lrcn --cnn_backbone resnet34 --pretrained true \
  --n_classes 50 --fr_per_vid 16 --batch_size 4 --mode train \
  --n_epochs 30 --output_dir outputs
\end{verbatim}

\begin{table}[h]
    \centering
    \caption{Training hyperparameters.}
    \label{tab:training-config}
    \begin{tabular}{ll}
        \toprule
        Hyperparameter & Value \\
        \midrule
        Optimizer & Adam \\
        Learning rate & $3 \times 10^{-5}$ \\
        Loss function & Cross-entropy \\
        Batch size & 4 \\
        Epochs & 30 \\
        Scheduler & ReduceLROnPlateau \\
        Device & \placeholder{insert CPU or GPU details} \\
        \bottomrule
    \end{tabular}
\end{table}

\section{Experimental Results}

\subsection{Training and Validation Progress}

Training logs are saved to \texttt{outputs/training_history.json}.
After the full run, insert a plot or table summarizing train and validation loss and accuracy across epochs.

\begin{figure}[h]
    \centering
    \fbox{\parbox[c][2.0in][c]{0.85\linewidth}{\centering \placeholder{Insert training/validation loss and accuracy plot here.}}}
    \caption{Training and validation progress over epochs.}
    \label{fig:training-curves}
\end{figure}

\begin{table}[h]
    \centering
    \caption{Final training summary. Replace placeholders after training.}
    \label{tab:training-summary}
    \begin{tabular}{lc}
        \toprule
        Metric & Value \\
        \midrule
        Best validation accuracy & \placeholder{insert value} \\
        Epoch of best checkpoint & \placeholder{insert value} \\
        Final training loss & \placeholder{insert value} \\
        Final validation loss & \placeholder{insert value} \\
        \bottomrule
    \end{tabular}
\end{table}

\subsection{Test Metrics}

Final test metrics are saved to \texttt{outputs/test_metrics.json}.
The assignment requires reporting multiclass classification metrics, including F1 score, AUC score, and a confusion matrix.
Table~\ref{tab:test-metrics} is reserved for the final test results.

\begin{table}[h]
    \centering
    \caption{Test-set metrics. Replace placeholders after evaluation.}
    \label{tab:test-metrics}
    \begin{tabular}{lc}
        \toprule
        Metric & Value \\
        \midrule
        Accuracy & \placeholder{insert value} \\
        Macro F1 & \placeholder{insert value} \\
        Weighted F1 & \placeholder{insert value} \\
        Macro AUC, one-vs-rest & \placeholder{insert value or N/A} \\
        \bottomrule
    \end{tabular}
\end{table}

\subsection{Confusion Matrix}

The confusion matrix is saved to \texttt{outputs/confusion_matrix.csv}.
After evaluation, insert a heatmap visualization of the confusion matrix and briefly discuss which classes are most frequently confused.

\begin{figure}[h]
    \centering
    \fbox{\parbox[c][2.25in][c]{0.85\linewidth}{\centering \placeholder{Insert UCF50 confusion matrix heatmap here.}}}
    \caption{Confusion matrix for the UCF50 test set.}
    \label{fig:confusion-matrix}
\end{figure}

\section{Discussion}

\placeholder{After the final experiment, discuss whether the model reached the required 65\% test accuracy threshold. Summarize likely causes of strong and weak performance, such as limited temporal context, visual similarity between action classes, class imbalance, or CPU/GPU training limits.}

\placeholder{Discuss any improvements attempted or planned, such as increasing frames per video, using a larger CNN backbone, increasing the number of epochs, tuning learning rate, using stronger augmentation, or adding a 3D CNN baseline.}

\section{Conclusion}

This project implements an LRCN pipeline for UCF50 action recognition, including preprocessing, stratified splitting, training, checkpointing, evaluation, and multiclass metric reporting.
The final version of this report should include the completed training curves, test metrics, confusion matrix, and a short analysis of whether the required classification accuracy was achieved.

\end{document}
